\section{Integral $K$-groups: concrete known consequences and the missing export}
\label{sec:source-K-bridge}

The final footnote in Alp\"oge's four-page remark suggests consequences for
$K$-groups of $\mathbb Z$, but the supplied declarations do not identify a
degree, a torsion prime, an eigenspace, or a comparison map.  There are
nevertheless precise existing integral results with which such an export
can be compared.  We keep ordinary integral $K$-groups, their reductions
modulo $p$, and their $p$-primary subgroups distinct.

\subsection{The exact cyclotomic comparison}

\paragraph{Finding F-SM-14 (the mod-$p$ comparison and its vanishing consequence).}
Let $p$ be an odd prime, let $n\geq1$, put
$F_p=\mathbb Q(\zeta_p)$, and let $C_p$ be the $p$-primary subgroup of
$\operatorname{Cl}(F_p)$.  Let $\omega$ be the Teichm\"uller character of
$\Delta_p=\operatorname{Gal}(F_p/\mathbb Q)$.  There is an isomorphism of
$\mathbb F_p$-vector spaces
\begin{equation}\label{eq:SM-Kclass}
 K_{2n}(\mathbb Z)/p
 \simeq H^2_{\mathrm{\acute et}}
       (\mathbb Z[1/p],\mu_p^{\otimes(n+1)})
 \simeq (C_p/pC_p)(\omega^{-n}).
\end{equation}
Here the last expression denotes a character eigenspace.  In particular,
\[
 K_{2n}(\mathbb Z)\{p\}=0
 \quad\Longleftrightarrow\quad C_p(\omega^{-n})=0.
\]
The comparison with \emph{mod-$p$} \'etale cohomology in the first
isomorphism is a known theorem, not a consequence of RH: see Weibel
(2013), Chapter~VI, Corollary~8.3, chapter page~48, following
Theorem~8.2, chapter page~47.  The following proof spells out its
class-group specialization and the vanishing equivalence.

\begin{proof}
Let $R=\mathbb Z[\zeta_p,1/p]$.  The map
$\operatorname{Spec}R\to\operatorname{Spec}\mathbb Z[1/p]$ is finite
\'etale Galois with group $\Delta_p$, of order $p-1$.  Taking invariants
on $\mathbb F_p[\Delta_p]$-modules is exact, by averaging; the
Hochschild--Serre spectral sequence therefore identifies cohomology over
$\mathbb Z[1/p]$ with the invariants of cohomology over $R$.

The Kummer sequence gives
\[
 0\longrightarrow\operatorname{Pic}(R)/p
 \longrightarrow H^2_{\mathrm{\acute et}}(R,\mu_p)
 \longrightarrow\operatorname{Br}(R)[p]\longrightarrow0.
\]
There is exactly one prime of $F_p$ above $p$, generated by
$1-\zeta_p$.  Thus localizing at $p$ leaves the ideal class group
unchanged.  The field $F_p$ is totally imaginary, and the local-invariant
sequence for the Brauer group of its ring of $S$-integers has only this
one finite place; the sum-of-invariants map is then injective.  Hence
$\operatorname{Br}(R)[p]=0$.  The Kummer identification is
$\Delta_p$-equivariant and reads
\[
 H^2_{\mathrm{\acute et}}(R,\mu_p)
       \simeq C_p/pC_p.
\]
Because $R$ contains $\mu_p$, the higher coefficient twist is constant
over $R$.  With its descent action retained, it gives
\[
 H^2_{\mathrm{\acute et}}(R,\mu_p^{\otimes(n+1)})
 \simeq (C_p/pC_p)\otimes_{\mathbb F_p}\mu_p^{\otimes n}.
\]
The last tensor factor has character $\omega^n$ modulo $p$.
Its invariant subspace is exactly the $\omega^{-n}$ eigenspace of
$C_p/pC_p$, proving~\eqref{eq:SM-Kclass} from the cited comparison.

The finite $p$-group $C_p$ splits into its Teichm\"uller eigenspaces,
since $p-1$ is invertible in $\mathbb Z_p$; reduction modulo $p$ commutes
with the corresponding idempotents.  A finite $p$-group $A$ has
$A/pA=0$ if and only if $A=0$: if $A=pA$, iteration gives
$A=p^rA=0$ for $p^r$ annihilating $A$.  Apply this first to each class
component, and then to the $p$-primary component of the finite group
$K_{2n}(\mathbb Z)$.  The latter finiteness is the usual integral
$K$-theory finiteness theorem, also used in the cited comparison.
\end{proof}

The assertion is a comparison of reductions modulo $p$ and hence a
vanishing criterion.  It does not identify the entire finite
$p$-primary groups with one another.  In positive-rank odd degrees the
correct comparison uses $p$-adic coefficients/completion; an
uncompleted finitely generated $K$-group cannot surject onto a copy of
$\mathbb Z_p$.  This distinction is material when reading informal
Chern-map summaries.

\subsection{A complete known vanishing and its class-group consequence}

Dutour Sikiri\'c, Elbaz-Vincent, Kupers, and Martinet (2026),
Theorem~6.4, prove
\[
 K_8(\mathbb Z)=0.
\]
The theorem is on page~19 of the author-posted updated manuscript;
Corollary~7.1 on page~20 concludes
$C_p(\omega^{p-5})=0$ for every prime $p>3$.  The journal article is in
\emph{Journal of the Institute of Mathematics of Jussieu},
\emph{25}(1), 565--590, DOI~10.1017/S1474748025101394;
its online publication date is 24 November 2025.  This is a known result
independent of the adopted Alp\"oge input.

The final deduction has a short, exact audit trail.  Their
Proposition~6.1 puts $H_9(BQ;\mathbb Z)$ in the class of finite groups
whose prime divisors are at most seven, and Proposition~6.3 puts the
kernel of the relevant Hurewicz map in the analogous class for five.
Since $K_8(\mathbb Z)=\pi_9(BQ)$, its possible prime divisors are at
most seven.  The odd primes $3,5,7$ are regular (their relevant Bernoulli
numerators are those of $B_2=1/6$ and $B_4=-1/30$), so Weibel VI,
Proposition~10.5, chapter page~63, removes their torsion.  The
$2$-primary part is zero by Weibel VI, Corollary~9.8, chapter page~55.
This last input makes explicit the prime-two step that is not written in
the updated manuscript's very short proof of Theorem~6.4.  The Voronoi
cell enumerations and the preceding homology computations are carried
as published inputs; they have not been rerun here.

Applying F-SM-14 with $n=4$ gives
$(C_p/pC_p)(\omega^{-4})=0$, hence
$C_p(\omega^{-4})=0$.  As $-4\equiv p-5\pmod{p-1}$, this is the stated
cyclotomic class-group consequence.  The argument works at the level of
finite quotients without an incorrect uncompleted Chern map.  It also
does not need the assertion, as printed in the manuscript's Section~7,
that $H^1_{\mathrm{\acute et}}(\mathbb Z[1/p],\mathbb Z_p(j))$
always vanishes when $p>j+1$.  For the positive twists relevant here,
that assertion requires even $j$: for odd $j>1$ this group has a
rank-one summand, as
Weibel VI, Exercise~8.3, chapter page~50, specifies.

\subsection{A density-one family with all degrees specified}

Chen et al. (2026), arXiv:2602.05090v1, Theorem~1.1 on page~2, prove
the following quantitative partial regularity statement.  For each
$\alpha>1/2$, let
\[
 M_\alpha(p)=\left\lfloor\frac{\sqrt p}{(\log p)^\alpha}\right\rfloor.
\]
The number of primes $p\leq X$ for which some even index
$2\leq2k\leq\min(M_\alpha(p),p-3)$ has
$p\mid\operatorname{num}(B_{2k})$ is
$O_\alpha(X/(\log X)^{2\alpha})$.  Their Corollary~1.4, page~3,
gives the corresponding vanishing of the even class-group components;
Corollary~1.7(4), page~5, gives the $K_{4k-2}$ consequence.
The formal verification reported by that paper concerns its main
elementary estimate; no Lean build of that archive has been performed
in this run, and the deep comparison inputs remain cited mathematics.

\paragraph{Finding F-SM-15 (both families of integral torsion consequences).}
For a fixed $\alpha>1/2$, outside
$O_\alpha(X/(\log X)^{2\alpha})$ odd primes $p\leq X$, the following
hold simultaneously for every integer $k$ with
$2\leq2k\leq\min(M_\alpha(p),p-3)$ and every integer $b\geq0$:
\begin{align*}
 K_{4k-2+2b(p-1)}(\mathbb Z)\{p\}&=0,\\
 K_{2p-2-4k+2b(p-1)}(\mathbb Z)\{p\}&=0.
\end{align*}
In particular, each fixed $\alpha$ gives a density-one set of primes
with these simultaneous vanishings.  The first line for $b=0$ is the
published consequence just cited; the second line makes explicit the
$4r$-degree consequence of its partial Vandiver components.

\begin{proof}
For a nonexceptional prime the partial regularity theorem gives
$p\nmid\operatorname{num}(B_{2k})$ throughout the range.  The
Herbrand--Ribet theorem identifies this with
$C_p(\omega^{1-2k})=0$; its exponent is also written $p-2k$ modulo
$p-1$.  Leopoldt reflection implies
$C_p(\omega^{2k})=0$.  These are exactly the two cited cyclotomic
inputs in Chen et al., equation~(4) and the deduction of
Corollary~1.4, page~3.

For the first line of the conclusion put
$n=2k-1+b(p-1)\geq1$ in F-SM-14.  Then
$-n\equiv1-2k\pmod{p-1}$.  For the second put
$n=p-1-2k+b(p-1)$.  The bound $2k\leq p-3$ ensures $n\geq2$, and
$-n\equiv2k\pmod{p-1}$.  Thus F-SM-14 proves both assertions, for
every $b$ at once.  They introduce no additional exceptional primes.
Finally, the prime number theorem gives
\[
 \frac{O_\alpha(X/(\log X)^{2\alpha})}{\pi(X)}
       =O_\alpha((\log X)^{1-2\alpha})\longrightarrow0.
\]
This proves the density assertion.
\end{proof}

The degrees in the second line are divisible by four because $p$ is
odd.  The degrees in the first line are congruent to two modulo four.
Both are precise $p$-primary assertions; neither line is an integral
vanishing assertion for an entire $K$-group at every prime.  Also, these
are pre-existing partial-regularity and comparison inputs, independent
of complex RH; their proof does not manufacture a missing exported
theorem from Alp\"oge's final footnote.

\subsection{The degree-one boundary and what an export must state}

\paragraph{Finding F-SM-16 (the elementary degree-one calculation).}
$K_1(\mathbb Z)\simeq\{\pm1\}\simeq\mathbb Z/2$.

\begin{proof}
Use the classical definition
$K_1(\mathbb Z)=\operatorname{GL}_\infty(\mathbb Z)^{\mathrm{ab}}$.
The determinant is a surjection onto $\{\pm1\}$.  Its kernel is
$\operatorname{SL}_\infty(\mathbb Z)$, which is generated by elementary
matrices: the Euclidean algorithm reduces the first column of a
unimodular matrix to a unit vector using integer row additions and signed
row interchanges; column additions clear the rest of its first row, and
induction applies to the remaining unimodular block.  Sign corrections
are allowed because
$\left(\begin{smallmatrix}0&1\\-1&0\end{smallmatrix}\right)
=e_{12}(1)e_{21}(-1)e_{12}(1)$ and its square is
$\operatorname{diag}(-1,-1)$.  In determinant one the final block also
has determinant one.  Stabilization supplies an extra index when needed.

For distinct indices $i,j,r$,
$e_{ij}(a)=[e_{ir}(a),e_{rj}(1)]$, so every elementary matrix is a
commutator in the stable group.  Conversely every commutator has
determinant one.  Thus the determinant kernel is precisely the
commutator subgroup, proving the asserted abelianization.
\end{proof}
The exact standard locator is Weibel III, Example~1.3.5, chapter
page~4, following Whitehead's Lemma~1.3.3.  This elementary nonzero
group already prevents a reading of the tentative footnote as blanket
vanishing in all positive degrees.

To make a further integral consequence of the supplied analytic method
reviewable, an export must specify its degree set, its coefficient or
torsion convention, its prime range, and the exact arithmetic map it
controls.  For example, an exported vanishing of the component
$C_p(\omega^{-n})$ would give the corresponding $p$-primary vanishing
in $K_{2n}(\mathbb Z)$ by F-SM-14.  To conclude an entire finite
$K_{2n}(\mathbb Z)$ is zero, one must cover every prime, including the
separate prime-two calculation.  Zero location for complex Dirichlet
$L$-functions alone supplies no such integral eigenspace or
divisibility statement in the supplied five declarations.

\subsection{Primary references and checked locators}

Chen, E., Cummins, C., Eltschig, B., Grubisic, D., Haller, L.,
Hong, L., Kurghinyan, A., Lau, K., Leather, H., Lee, S.,
Markosyan, A., Ono, K., Patel, M., Pendharkar, G., Rathi, V.,
Schneidman, A., Seeker, V., Sengupta, S., Sinha, I., \ldots\
Zhang, J. (2026).
\emph{Almost all primes are partially regular}
(arXiv:2602.05090v1).
\url{https://arxiv.org/abs/2602.05090}.
The used locators are Theorem~1.1, p.~2; equation~(4),
the reflection implication and Corollary~1.4, p.~3; and
Corollary~1.7(4), p.~5.  The natural-language proof of the main
quantitative theorem in Sections~2--3, pp.~6--10, was read throughout;
the formalization scope is stated in Section~4, pp.~10--11.

Dutour Sikiri\'c, M., Elbaz-Vincent, P., Kupers, A., \& Martinet, J.
(2026). Voronoi complexes in higher dimensions, cohomology of
$\operatorname{GL}_N(\mathbb Z)$ for $N\geq8$ and the triviality of
$K_8(\mathbb Z)$.
\emph{Journal of the Institute of Mathematics of Jussieu, 25}(1),
565--590.
\url{https://doi.org/10.1017/S1474748025101394}.
Updated author manuscript:
\url{https://jamartin.perso.math.cnrs.fr/Publications/DEKM.pdf}.
The used locators in this manuscript are Proposition~6.1, p.~18;
Proposition~6.3 and Theorem~6.4, p.~19; and Corollary~7.1, p.~20.
The complete final deductions were read; upstream Voronoi
computations were not reproduced.

Weibel, C. A. (2013).
\emph{The K-book: An introduction to algebraic K-theory}
(Graduate Studies in Mathematics, Vol.~145).
American Mathematical Society.
\url{https://doi.org/10.1090/gsm/145}.
Author's chapter files:
\url{https://sites.math.rutgers.edu/~weibel/Kbook/Kbook.III.pdf}
and
\url{https://sites.math.rutgers.edu/~weibel/Kbook/Kbook.VI.pdf}.
The exact locators are III, Example~1.3.5, chapter p.~4;
VI, equation~(8.1.1) and Theorem~8.2, chapter p.~47;
VI, Corollary~8.3, chapter p.~48; VI, Exercise~8.3, chapter p.~50;
VI, Corollary~9.8, chapter p.~55; VI, Proposition~10.5,
chapter p.~63; and VI, Theorem~10.9, chapter p.~65.
Chapter pagination differs from the printed volume; theorem
numbering is unchanged.  These primary statements and their
adjacent arguments were read directly.  The underlying motivic
comparison theorems are carried as established inputs.
